% docs/.paper/sec_intro.tex
\section{Introduction}
\label{sec:intro}

Software migration across programming languages is a frequent and consequential task in software engineering. Organisations regularly move legacy codebases written in C, C++, Java, or Go to modern systems-programming languages such as Rust. The transition removes critical memory-safety vulnerabilities (buffer overflows, use-after-free, data races)~\cite{szekeres2013sok,p58_syzygy,p56_rust_ffi_verification}, improves runtime efficiency, and lowers long-term maintenance cost. Manual repository migration, however, remains slow, expensive, and error-prone~\cite{khajeh2019automated,p46_muller_migration}.

Large Language Models (LLMs) translate isolated functions and synthesise unit tests with high accuracy~\cite{chen2021evaluating,roziere2020transcoder,p43_llm_testgen_survey}. Migrating an entire multi-file software repository introduces distinct technical barriers that defeat single-shot prompting. We identify four open challenges---language-pair engineering cost, hallucination under long horizons, the test-translation dichotomy, and transparent evaluation---each of which constrains the design of an autonomous translation harness.

\subsection{Four Challenges of Repository-Scale Translation}
\label{sec:intro_challenges}

\textbf{(1) Engineering cost per language pair.} Prior automated migration systems fall into two camps. Rule-based transpilers such as \textsc{C2Rust}~\cite{c2rust,p38_c2rust}, \textsc{CxGo}, and \textsc{Java2CSharp} use deterministic AST rewrites, but emit unidiomatic code with raw pointers and \texttt{unsafe} blocks~\cite{p50_rav1d_c2rust_case,p51_libexpat_c2rust_case} and require $\ge 10\,\text{K}$ lines of bespoke rewrites per source--target pair. Neuro-symbolic systems such as \textsc{AlphaTrans}~\cite{yang2024alphatrans,p02_alphatrans}, \textsc{Oxidizer}~\cite{siddiq2024oxidizer,p05_oxidizer}, and \textsc{Skel}~\cite{p64_skel} replace ad-hoc rules with program-analysis pipelines (\textsc{CodeQL}, type-driven proofs, execution traces), but their analysis components are likewise tied to one source--target pair and demand comparable engineering investment. Adding a second source language or a second target language typically requires rebuilding the analysis infrastructure from scratch.

\textbf{(2) Hallucination in long-horizon translation.} Repository-scale translation is a long-horizon task with feedback loops that span hundreds of model invocations. LLMs hallucinate type names, module paths, and trait bounds when context saturates or when the upstream callee has not yet been emitted~\cite{ibrahimzada2026recodeagent,p01_recodeagent,p02_alphatrans}. A monolithic prompt that contains the entire repository overwhelms small local models; a fragmented prompt that lacks the callee signatures of the module under translation causes the model to invent plausible-looking but non-existent symbols. Without an explicit ordering and bounded context, the system accumulates fabricated types that propagate through later repair passes.

\textbf{(3) The test-translation dichotomy.} Validating a translation requires executing test code in the target language. Translating tests is harder than translating source code~\cite{p45_translation_bugs} because test logic depends on the same type system as the source. A naive separation that places test translation and code translation in different agents fails when one buggy callee blocks every downstream test. Conversely, co-generating code and tests in a single agent biases the tests toward the (possibly wrong) generated code---the model weakens assertions to report a passing run rather than fix the underlying defect~\cite{p25_advtestgen,p02_alphatrans,p01_recodeagent}.

\subsection{The MAgHARCM Framework}
\label{sec:intro_magharcm}

We introduce \textbf{MAgHARCM} (\textbf{M}ulti-\textbf{Ag}ent \textbf{H}arness for \textbf{A}rchaeology-guided \textbf{R}epository-scale \textbf{C}ode \textbf{M}igration), an open-source translation harness written in Go that addresses all four challenges while running entirely on local Small Language Models (SLMs) served by Ollama. MAgHARCM coordinates four specialised agents over an explicit \textsc{Analyze} $\to$ \textsc{Plan} $\to$ \textsc{Translate} $\to$ \textsc{Validate} $\to$ \textsc{Repair} loop:

\begin{itemize}
    \item \textbf{Analyzer Agent} (\S\ref{sec:method_analyzer}) inspects the source AST and build manifest, then produces three documents: a Source Project Research report, a Third-Party Library Analysis mapping source libraries to idiomatic target crates, and a Target Project Design document that fixes a one-to-one structural mapping between source and target layouts.
    \item \textbf{Planning Agent} (\S\ref{sec:method_planning}) fragments the source AST into translation units, resolves ambiguous symbol names through an optional LSP-backed Navigator, builds a cycle-free reverse-topological dependency DAG, and emits a two-part \emph{Implementation Plan}---Part~A for source code translation and Part~B for test translation and validation.
    \item \textbf{Translator Agent} (\S\ref{sec:method_translator}) executes the plan in topological order, replacing skeleton stubs with concrete target code. Each fragment is synthesised with a bounded context ($\le 4$\,KB) assembled from a Symbol Navigator and a Prior-Modules Memory, keeping the local 4B coding model within its effective attention budget.
    \item \textbf{Validator Agent} (\S\ref{sec:method_validator}) executes the target compiler and test suite, emits a structured validation report, and feeds failures back to the Translator Agent for repair. The cascade combines AST syntax pre-checks, compiler build validation, test-suite execution, an adversarial test-weakening guard that halts the loop if assertions are removed, and a CodaMOSA-style coverage-plateau detector that prevents unproductive repair cycles.
\end{itemize}

A Strategy Registry (\S\ref{sec:method_strategies}) replaces a prior hard-coded cascade with a slice of named strategies matched against a repository profile; the analyzer iterates the registry and runs the first strategy that accepts the profile, falling through to subsequent strategies if the chosen one fails. This addresses challenge (1) by removing hard-coded behaviour and makes the harness extensible without editing the analyzer loop.

To address challenge (4), every iteration of the translate--validate loop writes a checkpoint to disk (\texttt{.artifacts/<run-id>/checkpoints/iter-NNNN.json}) and emits structured log events for every agent invocation. The empirical evaluation in \S\ref{sec:eval} reports per-iteration compilation status, per-iteration test pass rate, per-failure error codes, and a four-class root-cause failure taxonomy computed from those events.

\subsection{Key Contributions}
\label{sec:intro_contributions}

The primary contributions of this paper are:

\begin{enumerate}
    \item \textbf{Modular multi-agent architecture for local SLMs.} We formulate and implement MAgHARCM, a four-agent migration harness optimised for a 30B reasoning model paired with an instruction-tuned 4B coding model, both served by Ollama. The architecture addresses all four repository-translation challenges above without cloud inference.
    \item \textbf{Reverse-topological scheduling with bounded context.} We introduce cycle-free reverse-topological dependency scheduling with back-edge removal, paired with a $\le 4$\,KB Symbol Navigator and a $\le 4$\,KB Prior-Modules Memory that together cap the context per fragment and prevent saturation on small coding models.
    \item \textbf{Multi-stage validation and adversarial defense.} We design a Validator cascade that combines AST syntax pre-checks, compiler feedback, an adversarial test-weakening guard that halts the run if assertions are removed across iterations, and a CodaMOSA-style coverage-plateau detector that bounds local-SLM repair spinning.
    \item \textbf{Empirical evaluation and failure taxonomy.} We evaluate MAgHARCM on four open-source benchmark repositories across $K{=}3$ stochastic trials without artificial timeouts, report per-iteration compilation and test dynamics, and formalise a four-class root-cause failure taxonomy from 87 observed failures.
\end{enumerate}

The remainder of this paper is organised as follows. \S\ref{sec:method} formalises the problem, describes the tool surface, and details the four agents. \S\ref{sec:eval} presents the empirical evaluation. \S\ref{sec:related} surveys related work. \S\ref{sec:conclusion} concludes.
